\makeatother
\setcounter{counterchapter}{5}

\chapter{Effects of ice nucleating particles on modelled Arctic clouds}
\label{chap5}
\chaptermark{Effects of ice nucleating particles on modelled Arctic clouds}


 \epigraph{``The only way around is through.''}{\vspace{\baselineskip}Robert Frost}

%%%%%%%%%%%%%%%%%%%%%%%%%%%%%%%%%%%%%%%%%%%%%%

\newpage
\vspace*{\fill}
\begin{singlespace}
\textbf{Chapter content}

~

\startcontents[chapters]
\printcontents[chapters]{}{5}{}
\end{singlespace} 
\vspace*{\fill}
%%%%%%%%%%%%%%%%%%%%%%%%%%%%%%%%%%%%%%%%%%%%%%
\newpage

%%%%%%%%%%%%%%%%%%%%%%%%%%%%%%%%%%%%%%%%%%%%%%
\section{Introduction}
%%%%%%%%%%%%%%%%%%%%%%%%%%%%%%%%%%%%%%%%%%%%%%

Until now, this thesis tackled two thirds of the its announced plan through, the last two chapters. From the emission sources of INPs in \cref{chap3}, to the modelling of INP populations in the Arctic in \cref{chap4}.
This fifth chapter eventually addresses the question of the impact of particles responsible for ice nucleation on Arctic clouds, making a big step on the road toward the understanding of an hypothetical feedback loop that would link INPs and global warming regional impacts. 
We want to explore the possible effects of INPs on clouds by improving the representation of their interactions with the cloud microphysics in the model.
The study presented in this chapter happens to be the direct sequel of the latter, by introducing a new modelling set-up within WRF-Chem, inspired from the INP tracers emissions that enabled the modelling of Arctic INP emissions in \cref{chap4}.

Several studies have tackled the question of aerosols particles effect on cloud properties with simulations tools.
\citet{seifert_aerosol-cloud-precipitation_2012} have found that assumptions on INPs and CCN can have strong effect on cloud properties such as total condensate water content and glaciation of clouds,  and small effects was observed on the total averaged precipitation.
Later, \citet{fan_aerosol_2014} ran simulations to investigate the influence of CCN and INPs on precipitation in mixed-phase cloud. They studied two scenarios characterized by different cloud temperatures. They deduced various degrees of significance regarding the impacts of CCN and INPs on precipitation, but they clearly observed a more pronounced effect by INPs than CCN.
More recently, \citet{lee_simulations_2024} led investigations on INP and CCN effects on clouds with a focus on the ice crystal morphology. They observed that a higher level INP yields smaller ice crystal and irregular crystal shapes. Although crystals morphology is out of our focus in the present study, the size of ice crystals drives their number at constant water path; a measure that our modelling set-up is able to make.
In the Arctic, \citet{gjelsvik_using_2025} have demonstrated that a pure temperature-dependant ice nucleation parametrisation constructed on local observations can significantly improve the model representation of cloud radiative effect. Their investigations also showed that INPs are most likely important for accurate cloud representation in the Arctic, supporting the hypothesis that primary ice production is central for cloud phase.
Others studies have explored the role of specific particles\----especially dust\----in simulations of Arctic clouds for the investigations of the impact on their properties \citep{english_contributions_2014, shi_relative_2022, kawai_dominant_2023}, but literature lacks of studies that consider simultaneously the whole Arctic region and all the most important particles that can act as ice nuclei.
In this study, envision filling this gap.


%\citep{barthlott_aerosol_2018} --> only about ccn
%\todo{state of the science about INP in microphysics modelling}
\vspace{\baselineskip}
In a first section, will be presented and discussed the different microphysics parametrisations used in this study. It will be mainly the ice nucleation part of the schemes that will draw our focus. The modifications brought to the microphysics scheme of Thompson to link it to the WRF-Chem INPs will be reviewed in detail.
The simulations will set up over the well known period of the MOCCHA 2018 campaign, from August 1st to September 21st 2018.
Will be compared the different versions of the microphysics schemes (control and test) through their cloud properties outputs. We will get interested in the ice crystal number concentration, all the way to the cloud radiative effect.%, by way of the phase repartition.
Eventually, we will led the discussion about how INPs interact with regional climate, and how the latter may impact their emissions.


\paragraph{NB:} The work presented in this study remains preliminary. The results shown and the associated discussions reflect the state of our understanding and of the development of the study project that lies behind this chapter. Consequently, many lines of research remain unexplored but will be addressed in a future work, in continuation with this thesis. 

%%%%%%%%%%%%%%%%%%%%%%%%%%%%%%%%%%%%%%%%%%%%%%
\section{Microphysics scheme}
%%%%%%%%%%%%%%%%%%%%%%%%%%%%%%%%%%%%%%%%%%%%%%
    In the frame of this thesis, two microphysics parametrisations were used within WRF-Chem: the Morrison two-moment scheme \citep{morrison_impact_2009} and Thompson aerosol-aware \citep{thompson_study_2014}.
    Both of them provides mass mixing ratios and number concentrations of liquid and solid cloud particles.                                                      
    In this chapter, we focus on the Thompson aerosol-aware scheme. It is adapted in order to set up a comparison experiment that will highlight the effect of online modelling of aerosols particles on the modelled clouds.

% ____________________________________________
    \subsection{Linking aerosol particles to cloud microphysics}
    \label{sec:inp_to_mp}
    The microphysics scheme that receives our modifications is Thompson aerosol-aware \citep{thompson_study_2014}.
    % list the different type of nucleation and how they are modelled
    It handles ice nucleation through four distinct processes:
        immersion nucleation,
        deposition nucleation,
        homogeneous freezing of aqueous aerosols,
        and freezing of precipitating water drops.
    % we will modify only the immersion
    In the continuity of the precedent chapter where only immersion nucleation was studied, this experiment focus of the effect of aerosol-cloud interactions on the nucleation by immersion. Consequently, only the immersion parametrisation is modified.  

    The base set-up is identical to the one described in \cref{tab:wrf-chem_set-up}.
    The microphysics module is modified in order to run the intended experiment on INP effects. The parameters and schemes used are listed in \cref{tab:mp_INP_set-up}.

    \begin{table}[ht!]
        \centering
        \caption{WRF-Chem microphysics set-up of the control and test simulations.}
        \begin{tabular}{l|cc}
        \toprule
        Option name & Control & INP-mp \\
        \midrule 
            Microphysics scheme & \citet{thompson_study_2014} & \citet{thompson_study_2014} \\
            INP & \citet{cooper_ice_1986} & \textbf{INP tracers} \\
            Immersion nucleation & \citet{bigg_formation_1953} & \textbf{Explicit nucleation} \\
            Deposition nucleation & \citet{cooper_ice_1986} & \textbf{Explicit nucleation}\\
            Homogeneous nucleation & \citet{koop_water_2000} & \citet{koop_water_2000} \\
            Water friendly aerosols & Forced at \SI{500}{\per\cubic\cm} & Forced at \SI{500}{\per\cubic\cm} \\
        \bottomrule
        \end{tabular}
        \label{tab:mp_INP_set-up}
    \end{table}

    % how do we use the INP tracers
    The differences between the control and the test (named \emph{INP-mp}) set-up stand in the immersion and deposition modes of nucleation.
    In the original publication, \citet{thompson_study_2014} use the \citet{demott_predicting_2010} ice nucleation parametrisation (D10) with the model dust to predict the number of INPs. It is then used to influence the immersion nucleation that is handled by the \citet{bigg_formation_1953} scheme.
    In our control version, the number of INPs is computed with the temperature dependent scheme of \citet{cooper_ice_1986} (C86) rather than D10. 
    For deposition nucleation, the number of ice crystals ($n_i$) produced is directly determined by D10 in the original scheme \citep{thompson_study_2014} with aerosol-aware mode or by C86 otherwise. In our control set-up, we keep C86.

    In the INP-mp simulation, the INPs are driven by the number concentrations of the INP tracers of \cref{chap4}. 
    They are then used explicitly to produce ice crystals in both immersion and deposition modes. In order to prevent ice crystal formation by these modes outside of clouds, the cloud fraction is used to cap the ice formation:
    \begin{equation}
        n_i = N_{\text{INP}} * f_{\text{cloud}}
    \end{equation}
    where $f_{\text{cloud}}$ is the model cloud fraction.
    The INP tracers are emitted with immersion mode ice nucleation parametrisations, normally not suited to deposition nucleation.
    In order to take that into account, and still being able to use the INP tracers in deposition mode, the number of INPs capable of deposition nucleation is set to \perc{10} of the immersion INPs. This value relies on the observations by \citet{schrod_long-term_2020}. Even though it is not a consolidated value, it accounts for the likely difference between immersion and deposition. 

    The INP tracers are binned by temperature, with \tempc{5} steps, from \tempc{-5} to \tempc{-25}. They are emitted from aerosol particles emissions sources of relevant species for ice nucleation. The particles concentrations are converted into the INP tracers with adequate INAS based ice nucleation parametrisations.
    The aerosols species used for the INP tracers emissions are dust, sea sprays, fungal spores and bacteria from terrestrial vegetation.
    Once emitted, the tracers are subject to wet and dry removal, as well as in-cloud activation based on temperature conditions.
    The detailed production and transport mechanisms of the INP tracers are described in \cref{sec:method_inp_emission}.

    In the microphysics scheme, which INP tracer is used to feed $n_i$ depends on a temperature condition.
    The tracer of temperature $T_{\text{INP}}$ prescribes $n_i$ when $T \in [T_{\text{INP} - 5};\, T_{\text{INP}}]$.
    In the particular case of $T <$ \tempc{-25}, the INP25 tracer is used. And when $T>$ \tempc{-5}, $n_i$ is set to $10^{-6}$ \si{\per\liter}.
    For comparison, the default value of $n_i$ in the standard Thompson scheme (i.e. non-aerosol-aware) is \SI{1}{\per\liter}.
    
    % note about the CCN
    The aerosol engine used in the base WRF-Chem set-up, GOCART, misses some aerosols sources that have been identified as important sources for CCNs. \todo{what are they exactly? what reference?}
    Consequently, if the aerosol-aware microphysics scheme uses the aerosols produced by GOCART as water friendly aerosols (i.e. potential CCNs), it underestimates the formation of cloud water.
    A biased diagnostic of droplets number concentrations ensued, which affects the overall cloud representation in the simulations.
    In order to address this difficulty, as well as clearing any unwanted variability between the control and INP-mp simulations, the number of water friendly aerosols ($nwfa$) is prescribed and set to the fix value of \SI{500}{\per\cubic\cm}. \todo{Explain this value}

    

%% ____________________________________________
%    \subsubsection{Experimental modelling set-up}


%%%%%%%%%%%%%%%%%%%%%%%%%%%%%%%%%%%%%%%%%%%%%%
\section{Effects on cloud}
%%%%%%%%%%%%%%%%%%%%%%%%%%%%%%%%%%%%%%%%%%%%%%
In this section, will be reviewed the differences observed on cloud properties between the two microphysics set-ups: control and INP-mp. 
The comparison is performed on averaged values over the simulation period, which extends from August 1st to 21st September 2018.

%two possible cases:
%
%1. lot of effects --> for what reasons, what is the causal chain? are they good reasons? Are the results more realistic ?
%
%2. little effect --> is it because immersion has a small part? is it because the microphysics is not sensitive to INPs ?
%
%HYPOTHETICAL CAUSAL CHAIN: more INP implies less precipitations. We have less INPs in general so more precipitations, especially in mid-latitude. It dries the atmosphere which makes less available water vapour moving toward the Arctic. Less water vapour causes less clouds.
% dans tous les cas, même s'il y a un pb dans la simulation, on voit qu'il y a un effet en Arctic de la microphysique des INPs.

% ____________________________________________
    \subsection{Differences in ice nucleating particles concentrations}
    \label{sec:inp_diff}
    The main distinction between the control and the INP-mp simulations is contained within the method used to predict concentrations of ice friendly particles that can nucleate cloud ice. 
    The control set-up uses the C86 parametrisation that depends only on temperature. The INP-mp set-up, uses the INP tracers that are emitted and advected in the model, and conditions on temperature for their activations.
    Because the INP tracers are binned in temperature bins of \tempc{5}, they are not easily comparable to the INP concentrations produced by C86.
    In order de compare their prognostics on INP concentrations, the INP tracers concentrations are combined with a temperature dependence that matches how they are activated by the microphysics module (c.f. \cref{sec:inp_to_mp}).
    The concentrations are then averaged over the simulation period, and between the surface and \SI{4000}{\meter} altitude.
    The corresponding maps are shown on \cref{fig:INP_control_test}.

     \begin{figure}[htb]
        \centering
        \label{fig:INP_control_test}
        \includegraphics[width=1\textwidth]{./chap5/images/temperature_C86_INPtracers.png}
         \caption[Arctic temperature, and difference between INP concentrations from C86 and INP tracers.]{Comparison of the period average of concentrations activated INP predicted by \citet{cooper_ice_1986} parametrisation (C86) and by the INP-mp set-up. The concentrations are averaged from surface to \SI{4000}{\meter} altitude. The temperature maps highlights the temperature dependence of C86, and shows that activation temperatures are rarely reached.}
    \end{figure}

    The difference is striking: C86 produces several orders of magnitude more INPs than the INP-mp set-up with INP tracers.
    This is explained by the differences in the dependencies of the two approaches. While C86 only depends on temperature, the INP-mp set-up is sensitive to temperature conditions and to the advection of the INP tracers.
    Figure~\cref{fig:INP_control_test} perfectly illustrates the temperature dependence of C86, of which the map is a ``negative'' of the temperature map, as seen in \cref{chap4} on \cref{fig:INP_spectra}.
    This logarithmic dependence on temperature of these kind of ice nucleation parametrisations can cause high concentrations at places where particles able to nucleate ice have little reasons to be.
    On the INP C86 map, the highest concentrations are found over the Greenland ice cap. This is because the lowest temperatures are reached there. However, the Greenland ice cap contains no source of potential INPs, and by its altitude, is remote from known sources.  

    Conversely, the map of INP-mp shows less homogeneous concentrations and more spatial variability because of its dependence on INP tracers concentrations.
    This double constraint (temperature and INP tracers concentrations), causes the concentrations of the INP-mp set-up to be much lower than what C86 predicts.
    As seen in the previous chapters, INP concentrations in Arctic are low.
    The analysis of the performances of the INP tracers in \cref{chap4} showed that they reproduced the observed levels of INPs in central Arctic.
    %ce qui fait que ça peut être considéré comme fidèle.
    % si la température est également bonne, alors, c'est que la carte de INP-mp repose sur des fondements davantage physique que celle de Cooper 86.
    % il faut donc considérer que ces concentrations très basses d'INP activés sont vraisemblables.


% ____________________________________________
    \subsection{Investigating the impact on cloud properties}
    % --- Cloud cover
    The cloud property that reflects the most significant changes in the microphysics, is the cloud cover.
    On \cref{fig:modelled_cloud_comparison}A clearly appears the effects of the INP-mp modifications on the total cloud cover. 
    The difference of the mean cloud cover computed over the Arctic region ($\theta >60$°N) accounts for a reduction of \perc{5}.
    The difference lies mostly in the Arctic, especially over the areas covered by sea ice.
    Actually, the delimitation is very clear, suggesting a correlation with how INP concentrations are affected by sea ice.
    In the control simulation, INP concentrations do not rely on actual particle emissions, but only on temperature with the C86 parametrisation. Consequently, only the effects of sea ice on temperature can influence heterogeneous ice nucleation.
    In the INP-mp set-up, INPs are modelled particles transported from various emission sources that are not only affected by temperature. 
    Marine emissions are inhibited by the presence of ice at the surface of the ocean. Thereby, INPs get scarcer in average in the areas of icy ocean.
    This observation is directly linked to the results of \cref{sec:second_approach}, especially to the diagnostic of INPs spatial distribution shown on \cref{fig:INP_maps_MOCCHA}: at the temperatures encountered in summer, the INP concentrations in central Arctic are mostly under $10^{-1}$ \si{\per\liter} in average. 

    % greenland coasts signal
    Coasts of Greenland shows the biggest difference, with lower cloud cover by up to \perc{60}.
    The signal spatially correlates with the emissions sources of INPs detected in \cref{sec:artofmelt_sources} during the analysis of the ARTofMELT 2023 campaign.
    Indeed, the detected sources (c.f. \cref{fig:source_INP_artofmelt}) were distributed along the Greenland coasts. The analysis concluded on mineral dust emissions.
    However, the data used to update the emissions of this kind of particles did not consider these areas. Consequently, the modelled mineral dust lacks in this area, as depicted on \cref{fig:emission_maps}.
    Assuming that C86 better represents the INP concentrations over Greenland coasts, the missing mineral dust emissions are the causes for the ebbs of cloud cover observed in the INP-mp simulation.
    This supports the conclusion of the ARTofMELT INP sources analysis that Greenland coasts are important sources of INPs.

   % the signal over the eurasia
    There is also fewer cloud cover over continental areas, especially at mid-latitude over Eurasia.
    WHY???

    % --- IWP and LWP
    The analysis of ice water path (IWP) and liquid water path (LWP), expressed in \si{\gram\per\square\meter}, brings elements for the understanding of the changes observed in cloud cover.
    The IWP is compared on panel B of \cref{fig:modelled_cloud_comparison}.
    The figure shows a global increase of the ice content of clouds. Over the whole \[domain/Arctic\] area, the INP-mp set-up produces \perc{17.8} more ice.
    It results in a decrease in LWP of \perc{12.7} (\cref{fig:modelled_cloud_comparison}C).
    This LWP difference happens mostly in central Arctic, in areas similar to where the cloud cover difference happens. 

    % number concentration of droplets and ice crystals
    Looking at the number concentrations of liquid droplets and ice crystals, we notice decreases of respectively \perc{8.8} and \perc{7.3}.
    For variation is analogous to the evolution of water path for liquid droplets, which means an almost constant droplet size, but ice crystals have clearly grown in the INP-mp simulation.
    Considering that mass scales with diameter as $m \propto D^b$, with $b$ a scaling factor dependant on crystal shape, and that $b \approx 2$ after \citet{heymsfield_improved_2010}, the increase in ice crystal diameter expresses as,
    \begin{equation}
        \Delta D_{\text{ice}} = \left(\frac{1 + \Delta IWP}{1 + \Delta N_{\text{ice}}}\right)^{1/b}
    \end{equation}
    which gives $\Delta D = 1.127$, or a \perc{12.7} increase in ice crystal diameter.

     \begin{figure}[htbp]
        \centering
        \label{fig:modelled_cloud_comparison}
        \includegraphics[width=1\textwidth]{./chap5/images/.png}
         \caption[Comparison of cloud properties between the control and the INP-mp simulations.]{}
    \end{figure}

% ____________________________________
    \subsection{Sub-regional impact}

\begin{table}[htbp]
\centering
\caption{Regional mean values of cloud and precipitation variables from the control and modified (INP-mp) simulations, along with their relative differences across ATL, PAC, GRE, NNA, SIB, and NP.}
\label{tab:mean_values_3lines}
\resizebox{\textwidth}{!}{
\begin{tabular}{llcccccc}
\hline
\textbf{Variable} & \textbf{Simulation} & \textbf{ATL} & \textbf{PAC} & \textbf{GRE} & \textbf{NNA} & \textbf{SIB} & \textbf{NP} \\
\hline
CLT & Control & 0.88 & 0.87 & 0.66 & 0.67 & 0.65 & 0.93 \\
    & INP-mp & 0.81 & 0.70 & 0.50 & 0.57 & 0.59 & 0.60 \\
    & Rel. diff. (\%) & \textminus7.70 & \textminus20.26 & \textminus23.20 & \textminus14.69 & \textminus9.68 & \textminus35.94 \\
\hline
IWP (g/m$^2$) & Control & 3.25 & 2.84 & 3.56 & 4.47 & 3.44 & 2.89 \\
               & INP-mp & 5.02 & 4.33 & 4.06 & 6.92 & 5.30 & 3.06 \\
               & Rel. diff. (\%) & +54.66 & +52.36 & +13.88 & +54.60 & +54.19 & +5.94 \\
\hline
ISWP (g/m$^2$) & Control & 54.03 & 48.16 & 54.36 & 78.47 & 58.46 & 49.11 \\
                & INP-mp & 57.14 & 50.83 & 54.46 & 81.93 & 61.20 & 49.53 \\
                & Rel. diff. (\%) & +5.76 & +5.55 & +0.18 & +4.41 & +4.67 & +0.86 \\
\hline
LWP (g/m$^2$) & Control & 124.44 & 101.46 & 32.98 & 78.69 & 84.95 & 88.05 \\
               & INP-mp & 73.71 & 42.44 & 10.94 & 45.44 & 54.21 & 17.97 \\
               & Rel. diff. (\%) & \textminus40.77 & \textminus58.17 & \textminus66.81 & \textminus42.26 & \textminus36.18 & \textminus79.60 \\
\hline
VWP (g/m$^2$) & Control & 1.55$\times$10$^{4}$ & 1.40$\times$10$^{4}$ & 7.43$\times$10$^{3}$ & 1.37$\times$10$^{4}$ & 1.65$\times$10$^{4}$ & 1.14$\times$10$^{4}$ \\
               & INP-mp & 1.54$\times$10$^{4}$ & 1.39$\times$10$^{4}$ & 7.40$\times$10$^{3}$ & 1.36$\times$10$^{4}$ & 1.65$\times$10$^{4}$ & 1.14$\times$10$^{4}$ \\
               & Rel. diff. (\%) & \textminus0.78 & \textminus0.42 & \textminus0.40 & \textminus0.14 & \textminus0.29 & \textminus0.46 \\
\hline
$N_{ice}$ (m$^{-3}$) & Control & 1.74$\times$10$^{7}$ & 1.56$\times$10$^{7}$ & 2.09$\times$10$^{7}$ & 2.43$\times$10$^{7}$ & 1.86$\times$10$^{7}$ & 1.49$\times$10$^{7}$ \\
                     & INP-mp & 1.44$\times$10$^{7}$ & 1.36$\times$10$^{7}$ & 1.64$\times$10$^{7}$ & 2.17$\times$10$^{7}$ & 1.59$\times$10$^{7}$ & 1.01$\times$10$^{7}$ \\
                     & Rel. diff. (\%) & \textminus17.27 & \textminus12.81 & \textminus21.49 & \textminus10.88 & \textminus14.59 & \textminus32.07 \\
\hline
RAINNC (mm) & Control & 69.38 & 61.78 & 77.84 & 89.74 & 78.00 & 42.92 \\
             & INP-mp & 72.15 & 63.62 & 73.12 & 93.63 & 83.77 & 40.73 \\
             & Rel. diff. (\%) & +4.00 & +2.99 & \textminus6.05 & +4.34 & +7.39 & \textminus5.10 \\
\hline
\end{tabular}
}
\end{table} 
     \begin{figure}[p]
        \centering
        \label{fig:map_and_profiles}
        \includegraphics[width=1\textwidth]{./chap5/images/Arctic_map_profiles.png}
         \caption[Regional profiles of Arctic cloud properties.]{Ice water path, liquid water path, and cloud fraction vertical profiles from the control and the INP-mp simulations from six Arctic regions.}
    \end{figure}

% ____________________________________________
    \subsection{Comparison with lidar-radar data}

%%%%%%%%%%%%%%%%%%%%%%%%%%%%%%%%%%%%%%%%%%%%%%
\section{Effect of high latitude aerosols emission}
%%%%%%%%%%%%%%%%%%%%%%%%%%%%%%%%%%%%%%%%%%%%%%
In this section, we focus on the differences between two versions of the INP-mp simulation set-up.



%%%%%%%%%%%%%%%%%%%%%%%%%%%%%%%%%%%%%%%%%%%%%%
\section{Discussion about the possible feedback effects}
%%%%%%%%%%%%%%%%%%%%%%%%%%%%%%%%%%%%%%%%%%%%%%

What about the original question: is there a possible feedback loop between INPs and regional climate?

We have to determine if we have more or less INPs, and where.

    \vspace{\baselineskip}
    \todo{Fill the Key results box}
    \boxKR{Modelled Arctic INPs and related effects on clouds}{
        \begin{itemize}
            \item 
        \end{itemize}
    }
%%%%%%%%%%%%%%%%%%%%%%%%%%%%%%%%%%%%%%%%%%%%%%
\stopcontents[chapters]
